\chapter{Triển khai và vận hành}

\section{Công nghệ và tổ chức mã nguồn}

Backend được viết bằng Python 3.11. Package \texttt{rbac\_guard} tách parser,
rules, scoring, reporting, context profiler, repository RBAC, application
service, CLI và FastAPI adapter thành các module có trách nhiệm riêng. SQLite
nằm trong thư viện chuẩn; pytest và coverage kiểm thử lõi. Streamlit là giao
diện tùy chọn để tải log, lọc alert và tải kết quả.

Demo Nova Bank dùng FastAPI làm API thật và Next.js 15/React 19 làm giao diện.
Frontend không lưu business rule hoặc danh sách quyền cứng. Nó gọi API để lấy
phiên, danh sách role, user, account, giao dịch, policy và audit log; backend
là nơi duy nhất quyết định kết quả authorization. Cách tổ chức này giúp cùng
một SQLite database thể hiện được thay đổi tài khoản, role, policy và giao dịch
trong buổi trình bày.

\section{Các giao diện vận hành}

\begin{table}[htbp]
\centering
\small
\begin{tabularx}{\textwidth}{p{0.28\textwidth}X}
\toprule
Giao diện & Công dụng \\
\midrule
\texttt{rbac-guard init-db} & Tạo SQLite schema và nạp seed RBAC có kiểm tra tham chiếu. \\
\texttt{rbac-guard check-access} & Trả lời một user có quyền resource:action hay không. \\
\texttt{rbac-guard analyze} & Parse log, phát hiện, chấm điểm và ghi artifact; nhận tùy chọn \texttt{--context-risk}. \\
\texttt{rbac-guard evaluate} & Đối chiếu event ID với expected label; tùy chọn manifest để tính thêm metric campaign. \\
\texttt{rbac-guard export-audit-events} & Chuyển audit Nova Bank thành CSV Event cho pipeline. \\
\texttt{GET /health} & Kiểm tra FastAPI đang lắng nghe và database được chọn. \\
\texttt{/auth/*}, \texttt{/users}, \texttt{/transactions} & Cung cấp xác thực, quản trị nhân viên và nghiệp vụ giao dịch có kiểm soát RBAC. \\
\path{/approvals/{reference}} & Bề mặt kiểm chứng trực tiếp separation of duties. \\
\bottomrule
\end{tabularx}
\caption{Các giao diện quan trọng của hệ thống.}
\label{tab:interfaces}
\end{table}

\section{Thiết lập}

Môi trường Python được đồng bộ bằng \texttt{uv}; frontend cài dependency trong
thư mục \texttt{web}. Backend demo chạy trên cổng 8000, frontend chạy trên cổng
3000. CORS chỉ cho origin frontend cấu hình bởi \texttt{RBAC\_WEB\_ORIGIN};
đường dẫn SQLite có thể thay đổi bằng \texttt{RBAC\_DEMO\_DB}. Các biến môi
trường này giúp chạy lại demo mà không phải sửa mã nguồn.

\begin{lstlisting}[language=bash]
uv sync --extra api
uv run uvicorn rbac_guard.api:app --host 127.0.0.1 --port 8000
cd web && npm install && npm run dev
\end{lstlisting}

Quy trình đánh giá holdout được tự động hóa bởi
\path{scripts/generate_holdout_dataset.py} và
\path{scripts/run_holdout_sensitivity.py}. Script sinh dữ liệu dùng thứ tự
và giá trị cố định, ghi manifest cùng phân tầng; script độ nhạy chạy lưới ngưỡng
3/5/7 và cửa sổ 60/300/600 giây. Hai artifact metric được giữ trong repository
để bảng kết quả có thể đối chiếu với đầu ra máy thay vì số nhập tay.

\section{Kiểm thử đường tích hợp Nova Bank--pipeline}

Integration test khởi tạo database, tạo năm lần login thất bại và một lần
Teller truy cập quyền phê duyệt bị từ chối. Adapter xuất toàn bộ audit theo thứ
tự tăng dần; parser đọc lại file; application service phân tích và phải tạo cả
\texttt{password\_guessing} lẫn \texttt{unauthorized\_access}. Test CLI riêng
kiểm tra lệnh export tạo file có thể parse. Đây là bằng chứng tích hợp tối thiểu
giữa hai phần của dự án; SQL Injection vẫn cần log HTTP vì audit Nova Bank
không giữ raw request.

\section{Kịch bản demo Nova Bank}

Kịch bản 5--7 phút bắt đầu bằng Admin tạo tài khoản Teller. Teller đăng nhập,
tạo yêu cầu chuyển khoản 50.000.000 VND và hoàn tất giao dịch trong policy
baseline; đây là đối chứng cho rủi ro một người tạo và duyệt. Admin sau đó đổi
sang \emph{Phân tách nhiệm vụ}. Teller tạo giao dịch thứ hai rồi mở trực tiếp
URL phê duyệt, nhưng backend từ chối. Cuối cùng Controller đăng nhập, phê duyệt
giao dịch hợp lệ và Admin kiểm tra audit log. Hướng dẫn chi tiết, dữ liệu mẫu
và thao tác khôi phục trạng thái được giữ trong \texttt{DEMO\_RBAC\_NOVA\_BANK.md}.

Việc mở trực tiếp URL là bước kiểm chứng quan trọng: nó phân biệt enforcement
ở backend với một giao diện chỉ ẩn nút. Trạng thái giao dịch phải tiếp tục là
pending sau lần từ chối, trước khi Controller thực hiện phê duyệt.
